\documentclass[11pt,a4paper]{article}
\usepackage[utf8]{inputenc}
\usepackage[T1]{fontenc}
\usepackage[english]{babel}
\usepackage{geometry}
\geometry{margin=2.5cm}
\usepackage{setspace}
\setstretch{1.08}
\usepackage{hyperref}
\usepackage{longtable}
\usepackage{booktabs}
\usepackage{listings}
\usepackage{xcolor}
\usepackage{enumitem}
\hypersetup{colorlinks=true,linkcolor=blue,urlcolor=blue}

\lstdefinestyle{cmd}{
  basicstyle=\ttfamily\small,
  frame=single,
  breaklines=true,
  showstringspaces=false
}

\title{Meso2D: User Tutorial}
\author{Practical guide for new users}
\date{\today}

\begin{document}
\maketitle
\tableofcontents
\newpage

\section{Purpose of this tutorial}
This guide explains how to use Meso2D to generate 2D mesostructures with:
\begin{itemize}[leftmargin=1.5em]
  \item aggregates (circular, elliptical, or polygonal),
  \item pores,
  \item reactive points,
  \item and exportable outputs for post-processing.
\end{itemize}

The approach is practical: follow a clear sequence to run reproducible simulations and export consistent data.

\section{What the program does}
Meso2D builds a 2D specimen and fills it with geometric entities according to:
\begin{itemize}[leftmargin=1.5em]
  \item gradation/dosification,
  \item aggregate area fraction,
  \item pore area fraction,
  \item reactive-point area fraction,
  \item and random seed control.
\end{itemize}

The execution is orchestrated by a simulation controller that selects a worker for:
\begin{itemize}[leftmargin=1.5em]
  \item \textbf{circulos},
  \item \textbf{ellipses},
  \item \textbf{poligonos}.
\end{itemize}

\section{Prerequisites}
\subsection{Dependencies}
Make sure at least the following packages are installed:
\begin{itemize}[leftmargin=1.5em]
  \item Python 3.x
  \item PyQt5
  \item numpy
  \item matplotlib
  \item ezdxf
  \item shapely
  \item opensimplex
\end{itemize}

\subsection{Run commands}
From the project root:
\begin{lstlisting}[style=cmd]
python Meso2D_refactor.py
\end{lstlisting}

For the English-oriented branch:
\begin{lstlisting}[style=cmd]
python Meso2D_refactor_EN.py
\end{lstlisting}

\section{Quick workflow (5 steps)}
\begin{enumerate}[leftmargin=1.7em]
  \item Choose the geometry mode (circles, ellipses, or polygons).
  \item Set specimen size and gradation parameters.
  \item Optionally enable pores and/or reactive points.
  \item Run the simulation and inspect visualization.
  \item Export structure and image.
\end{enumerate}

\section{Start screen: mode selection}
At startup, the mode selection window determines:
\begin{itemize}[leftmargin=1.5em]
  \item aggregate placement algorithm,
  \item geometric representation of results,
  \item collision-check strategy.
\end{itemize}

\textbf{Recommendation:}
\begin{itemize}[leftmargin=1.5em]
  \item Start with \textit{circulos} for fast checks.
  \item Use \textit{ellipses} or \textit{poligonos} for more realistic morphology studies.
\end{itemize}

\section{Project configuration}
\subsection{Main parameters}
\begin{longtable}{p{0.30\textwidth} p{0.64\textwidth}}
\toprule
\textbf{Parameter} & \textbf{Meaning and effect} \\
\midrule
x, y & Specimen dimensions (width and height). Total area is A = x*y. \\
Pagg & Target area fraction occupied by aggregates. \\
Gradation table sieve\_size / tpp & Defines the gradation curve. Area is split by fraction using these columns. \\
dporo\_min, dporo\_max & Minimum and maximum pore diameters. \\
Pporos & Target pore area fraction. \\
dpto\_min\_aggregates, dpto\_max\_aggregates & Diameter range for reactive points on aggregates. \\
Ppto\_react\_aggregates & Reactive-point area fraction over aggregates. \\
dpto\_min\_pasta, dpto\_max\_pasta & Diameter range for reactive points in paste. \\
Ppto\_react\_pasta & Reactive-point area fraction in paste. \\
seed & Random seed for exact reproducibility. \\
\bottomrule
\end{longtable}

\subsection{Gradation table}
\begin{itemize}[leftmargin=1.5em]
  \item Use \textbf{Add} to insert rows.
  \item Fill \textbf{sieve\_size} and \textbf{tpp} consistently.
  \item Physically consistent values improve convergence and placement success.
\end{itemize}

\section{Simulation options}
\subsection{Pores}
When pores are enabled:
\begin{itemize}[leftmargin=1.5em]
  \item pore radii are generated until target pore area is consumed,
  \item pores are placed without overlap against aggregates or existing pores.
\end{itemize}

\subsection{Reactive points}
Reactive points can be generated in two channels:
\begin{itemize}[leftmargin=1.5em]
  \item \textbf{on aggregates},
  \item \textbf{in paste}.
\end{itemize}

Internally:
\begin{itemize}[leftmargin=1.5em]
  \item aggregate channel typically uses \textbf{cluster sampling},
  \item paste channel typically uses \textbf{Simplex noise gating}.
\end{itemize}

\section{Running the simulation}
\subsection{Before pressing Run}
Recommended checklist:
\begin{enumerate}[leftmargin=1.7em]
  \item Correct mode selected.
  \item Valid x and y values.
  \item Complete gradation table.
  \item Pagg set.
  \item If pores/reactive points are enabled, all required fields are filled.
  \item Fixed seed if strict reproducibility is needed.
\end{enumerate}

\subsection{During execution}
\begin{itemize}[leftmargin=1.5em]
  \item Progress bar advances by pipeline stages.
  \item Information panel shows worker status messages.
  \item \textbf{Stop} can cancel the current run.
\end{itemize}

\subsection{After execution}
After a successful run:
\begin{itemize}[leftmargin=1.5em]
  \item structure visualization is available,
  \item export actions are enabled,
  \item simulation time is logged.
\end{itemize}

\section{Visualization and image saving}
\begin{itemize}[leftmargin=1.5em]
  \item \textbf{View structure}: opens the viewer window.
  \item \textbf{Save Image}: saves PNG/JPEG/TIFF/BMP.
\end{itemize}

\section{Structure export}
\subsection{Available formats}
From \textbf{Export Structure}, you can save to:
\begin{itemize}[leftmargin=1.5em]
  \item DXF
  \item SVG
  \item JSON
  \item Gmsh GEO
\end{itemize}

\subsection{Best practices}
\begin{itemize}[leftmargin=1.5em]
  \item Export right after a successful run.
  \item If parameters change, rerun before exporting.
  \item Include date and seed in file names for traceability.
\end{itemize}

\section{Project lifecycle actions}
\begin{itemize}[leftmargin=1.5em]
  \item \textbf{New Project}: clears current state.
  \item \textbf{Open Project}: loads a saved project and repopulates UI fields.
  \item \textbf{Save Project}: stores current parameters.
  \item \textbf{Load Example}: loads a sample project for quick testing.
\end{itemize}

\section{Reproducibility and traceability}
To reproduce exactly the same structure:
\begin{enumerate}[leftmargin=1.7em]
  \item save the project,
  \item record the seed,
  \item keep the same mode,
  \item run the same code version.
\end{enumerate}

\section{Troubleshooting}
\subsection{Simulation does not start}
Common causes:
\begin{itemize}[leftmargin=1.5em]
  \item missing required field,
  \item incomplete gradation table,
  \item non-numeric input in numeric fields.
\end{itemize}

\subsection{Too many placement failures}
\begin{itemize}[leftmargin=1.5em]
  \item reduce target fractions (Pagg, Pporos, reactive ratios),
  \item increase specimen size (x, y),
  \item reduce minimum diameters,
  \item begin with circles mode for calibration.
\end{itemize}

\subsection{Export fails}
\begin{itemize}[leftmargin=1.5em]
  \item ensure a fresh structure has been generated,
  \item check write permissions in destination folder,
  \item test DXF export first as a baseline.
\end{itemize}

\section{Recommended usage habits}
\begin{enumerate}[leftmargin=1.7em]
  \item Keep project versions and seeds documented.
  \item Change one parameter group at a time.
  \item Save both image and structure for every useful run.
  \item Record mode, seed, Pagg, Pporos, and reactive ratios.
  \item For batch studies, calibrate in circles mode and then switch to ellipses/polygons.
\end{enumerate}

\section{Quick start example}
\begin{enumerate}[leftmargin=1.7em]
  \item Run \texttt{python Meso2D\_refactor\_EN.py}.
  \item Select \textit{circulos}.
  \item Set x=100, y=100 (example), Pagg, and gradation table values.
  \item Enable pores with moderate values.
  \item Press \textbf{Run}.
  \item Visualize and export to DXF + PNG.
\end{enumerate}

\section{Summary}
Meso2D combines gradation-driven area budgeting, stochastic sampling, and geometric collision constraints to build 2D mesostructures.

A robust routine is:
\begin{enumerate}[leftmargin=1.7em]
  \item configure parameters,
  \item run,
  \item validate visually,
  \item export,
  \item and record seed plus project file.
\end{enumerate}

This process provides reproducible and comparable outputs for later analysis.

\end{document}
